\labsection{6}{Sales Forecasting and S\&OP with ERP Data}{sec:lab06}

\begin{labproject}{Use the supplied planning spreadsheet and connect its inputs back to Odoo transactions}
\textbf{Coverage.} Tutorial 6 + Chapter 6.\
\textbf{Core systems.} Supplied Excel spreadsheet + Odoo Sales/Inventory/Manufacturing reporting.\
\textbf{Goal.} Build formula-driven forecast and S\&OP logic, then explain where an integrated ERP would supply the historical data used by the model.

\begin{labmode}
The tutorial explicitly requires the supplied spreadsheet. Odoo remains part of the lab by providing the operational context: Sales generates history, Inventory provides starting stock, and Manufacturing provides actual production. The spreadsheet performs the aggregate planning calculation.
\end{labmode}

\medskip\noindent\textbf{Part A --- Complete the simple sales forecast.}

Use the provided previous-year values. For each month:
\begin{enumerate}
  \item Previous Year Base = Previous Year Sales $-$ Previous Promotion;
  \item Growth = Previous Year Base $\times 3\%$;
  \item Base Projection = Previous Year Base + Growth;
  \item Forecast = Base Projection + Current Promotion.
\end{enumerate}

Current promotion assumptions: +500 pairs in November and +400 pairs in December.

Use formulas. Do not type the calculated rows manually.

\medskip\noindent\textbf{Part B --- Resolve units before capacity calculation.}

The tutorial states demand in pairs but capacity as 40 shoes/hour. Record which interpretation your tutor intends. Your formulas must use one consistent unit throughout.

\medskip\noindent\textbf{Part C --- Build the S\&OP plan.}

Calculate capacity from production rate, 8 hours/day, and monthly working days. Keep utilisation at or below 99\% as required in Tutorial 6. Use the inventory balance:
\[
I_t=I_{t-1}+P_t-F_t.
\]
Plan around the expected festive-demand pressure in June and December rather than merely copying forecast into production.

\medskip\noindent\textbf{Part D --- Check the source discrepancy.}

Tutorial 6 states safety stock = 100, while the supplied spreadsheet contains 500 in the initial December inventory cell. Do not hide the conflict. Confirm which value should be used for assessment and state it in your sheet or submission notes.

\medskip\noindent\textbf{Part E --- Inspect related Odoo reports.}

After completing earlier labs, open:
\begin{itemize}
  \item \pathmenu{Sales \textrightarrow Reporting};
  \item \pathmenu{Inventory \textrightarrow Reporting};
  \item \pathmenu{Manufacturing \textrightarrow Reporting}.
\end{itemize}

For each report, identify which S\&OP input or validation measure it could provide in a larger implementation: sales history, starting inventory, production output, or actual-versus-plan comparison.

\begin{keyidea}
Forecasting is not ``outside ERP.'' Even when the planning calculation is performed in a spreadsheet, the reliability of the model depends on integrated historical sales, inventory, and production data.
\end{keyidea}
\end{labproject}

\begin{labdeliverable}
Submit the screenshot of the completed forecast and S\&OP tables required by Tutorial 6. Include formulas, clearly state the unit interpretation and starting-inventory assumption, and add a short note naming the Odoo reports that would supply/validate the underlying operational data.
\end{labdeliverable}
